\documentclass[10pt,a4paper]{amsart}
\usepackage[margin=19mm]{geometry}
\usepackage{amsmath,amssymb,mathtools}
\usepackage[scr=rsfs,cal=euler]{mathalfa}
\usepackage{xfrac}
\usepackage[unicode,hidelinks,bookmarks=false]{hyperref}
\newcommand{\ie}{i.e.\ }
\newcommand{\cf}{cf.\ }
\newcommand{\resp}{respectively\ }
\newcommand{\Subsection}{\subsection}
\providecommand{\nobreakdash}{\nobreak\hbox{-}}
\setlength{\emergencystretch}{2em}
\multlinegap0pt
\usepackage{amssymb}
\usepackage{enumerate}
\newtheorem{lemma}{Lemma}
\newcommand{\BNorm}[2]{\Big\|#1\Big\|_{#2}}
\newcommand{\Norm}[2]{\|#1\|_{#2}}
\title{Weak-type characterizations of Sobolev and bounded variation spaces on metric measure spaces}
\author{Tuomas P.\ Hyt\"onen$^*$, Dachun Yang, Wen Yuan, Yirui Zhao}
\date{Source: arXiv:2608.24106v1 (CC BY 4.0)}
\begin{document}
\maketitle

\noindent\emph{Source and licence.} Weak-type characterizations of Sobolev and bounded variation spaces on metric measure spaces. Version arXiv:2608.24106v1 (CC BY 4.0), distributed by arXiv under the Creative Commons Attribution 4.0 International licence, http://creativecommons.org/licenses/by/4.0/, as recorded in the arXiv licence metadata for this exact version and shown on the abstract page https://arxiv.org/abs/2608.24106v1. Full licence notice: \texttt{licenses/arxiv-2608.24106-CC-BY-4.0.txt}.\par\smallskip

\noindent\textit{Editorial scope.} Counted unit: the article's lemma lem:BSVY105 (source0.tex lines 3196--3207). The article's complete original proof of it is delivered with it (source0.tex lines 3209--3368, one \texttt{proof} environment of 160 lines). No mathematics is written, repaired or re-derived: the statement and the proof are the article's own lines, in the article's own notation, and no retained line is substituted or re-flowed.  Retained uncounted as the source-local context the proof needs: lines 2960--2964; lines 3002--3012; lines 3018--3021.  Nothing was omitted from any retained span, so the package carries no omission record.  No retained line contains a citation, so the wrapper carries no bibliography and no bibliography entry.  No retained line contains a figure environment, an image-inclusion command or drawing code, so no image is delivered and the package declares no asset dependency.  Editorial formatting only, in addition: an AMS \texttt{amsart} preamble that restates the article's own declarations and macros used by the retained lines, namely the environments \texttt{lemma} on separate counters, together with the standard packages those lines need.  The retained environments print in this document as Lemma 1 in order of appearance, whereas the article prints the same items with the numbering of its own full document.  The complete native source of the article as distributed in the arXiv e-print for this exact version is bundled unchanged as \texttt{sources/208450/source0.tex}.
\begin{equation}\label{eq:ele}
\begin{split}
\int_{\{y\in X:[V(x,y)]^{-\alpha}>R\}}[V(x,y)]^{\alpha-1}\,d\mu(y)\lesssim R^{-1},
\end{split}
\end{equation}

\begin{lemma}\label{lem:BSVY963}
Let $(X,\rho,\mu)$ be a doubling metric measure space.
Let $p\in[1,\infty)$, $\gamma\in\mathbb{R}$,
$F\in L^\infty(\mu\times\mu)$, and $E\subset X$ be a bounded set. Then
\begin{equation*}
  \limsup_{s\nearrow 1}(1-s)^{\frac1p}\Norm{\mathbf{1}_{E\times
  E}F}{L^{p}(\rho^{-ps}V^{-1})}
  \lesssim\Norm{\rho^{-1-\gamma}F}{L^{p,\infty}(\rho^{p\gamma}V^{-1})},
\end{equation*}
where the implicit positive constant depends only on $p$ and $\gamma$.
\end{lemma}

\begin{equation}\label{eq:BSVY4-2}
  \int GH\,d\nu
  \leq q'\Norm{G}{L^{q,\infty}(\nu)}\Norm{H}{L^{q',1}(\nu)},\qquad q\in(1,\infty).
\end{equation}

\begin{lemma}\label{lem:BSVY105}
Let $(X,\rho,\mu)$ be a doubling metric measure space.
Let $p\in[1,\infty)$, $\gamma\in\mathbb{R}$,
$F\in L^\infty(\mu\times\mu)$, and $E\subset X$ be a bounded measurable set. Then
\begin{equation*}
  \limsup_{s\nearrow 1}(1-s)^{\frac1p}\Norm{\mathbf{1}_{E\times
  E}F}{L^{p}(\rho^{-ps}V^{-1})}
  \lesssim\Norm{\rho^{-1}V^{-\gamma}F}{L^{p,\infty}(V^{p\gamma-1})},
\end{equation*}
where the implicit positive constant depends only on $p,\gamma$, and the doubling constant
of $X$.
\end{lemma}

\begin{proof}
The proof of the present lemma is a technical modification of Lemma \ref{lem:BSVY963}. For
completeness, we give the details. As before, we consider different cases of
$\gamma\in\mathbb{R}$.

\emph{Case 0: $\gamma=0$}. For $\gamma=0$, the statement of Lemma \ref{lem:BSVY105} reduces
to the corresponding case of Lemma \ref{lem:BSVY963}, which we have already proved.

\emph{Case 1: $\gamma\in(0,\infty)$}. Let the symbol $D$ denote the upper dimension
of $X$, i.e., a positive constant such that, for all $x\in X$ and $0<r\le R<\infty$,
\begin{align}\label{eq:db}
\frac{V(x,R)}{V(x,r)}\lesssim \Big(\frac{R}{r}\Big)^D
\end{align}
with the implicit positive constant independent of $x,r$, and $R$.
Given $s\in(0,1)$, using the Lorentz space duality \eqref{eq:BSVY4-2} with $q:=
\frac{1+\gamma D}{s+\gamma D}\in (1,\infty)$, we obtain
\begin{equation}\label{eq:BSVY-v2}
\begin{split}
&\iint_{E\times E}\frac{|F|^p}{\rho^{ps}V}%\\
=\iint_{E\times E}\Big(\frac{|F|}{\rho^{s}V^\gamma}\Big)^pV^{p\gamma-1}    \\
&\quad=\iint_{X\times X}\Big(\frac{|F|}{\rho V^\gamma}\Big)^{\frac{p(s+D\gamma)}{1+D\gamma}}
\mathbf{1}_{E\times E}\Big(\frac{|F|
\rho^{D\gamma}}{V^\gamma}\Big)^{\frac{p(1-s)}{1+D\gamma}}
V^{p\gamma-1}\\
&\quad\le \frac{1+D\gamma}{1-s}
\BNorm{\Big(\frac{|F|}{\rho V^\gamma}\Big)^{\frac{p(s+D\gamma)}{1+D\gamma}}}{
L^{\frac{1+D\gamma}{s+D\gamma},\infty} (V^{p\gamma-1})} %\\
\BNorm{\mathbf{1}_{E\times E}\Big(\frac{|F|
\rho^{D\gamma}}{V^\gamma}\Big)^{\frac{p(1-s)}{1+D\gamma}}}{
L^{\frac{1+D\gamma}{1-s},1} (V^{p\gamma-1})}\\
&\quad\le \frac{1+D\gamma}{1-s}
\BNorm{\frac{F}{\rho V^\gamma}}{
L^{p,\infty}
(V^{p\gamma-1})}^{\frac{p(s+D\gamma)}{1+D\gamma}}\Norm{F}{L^\infty(\mu\times\mu)}
^{\frac{p(1-s)}{1+D\gamma}} %\\
\BNorm{\mathbf{1}_{E\times
E}\Big(\frac{\rho^{D\gamma}}{V^\gamma}\Big)^{\frac{p(1-s)}{1+D\gamma}}}{
L^{\frac{1+D\gamma}{1-s},1} (V^{p\gamma-1})}.
\end{split}
\end{equation}
Now, we estimate the last $L^{\frac{1+D\gamma}{1-s},1}$-norm.
Since $E$ is bounded, we infer that $E\subset B(x_0,R)$ for some $x_0\in X$
and $R\in(0,\infty)$. Clearly, for any $x,y\in E$, $B(x,\rho(x,y))\subset
B(x_0,3R)$. Applying this and \eqref{eq:db}, we find that, for any $x,y\in E$,
\begin{align}\label{eq:pvm}
\frac{\rho(x,y)^{D\gamma}}{V(x,y)^\gamma}\lesssim
\frac{R^{D\gamma}}{V(x_0,R)^\gamma}=:M.
\end{align}
Then we obtain
\begin{equation}\label{eq:lot1}
\begin{split}
&\BNorm{\mathbf{1}_{E\times
E}\Big(\frac{\rho^{D\gamma}}{V^\gamma}\Big)^{\frac{p(1-s)}{1+D\gamma}}}{
L^{\frac{1+D\gamma}{1-s},1} (V^{p\gamma-1})}\\
&\quad=\int_{0}^{\infty}\Big[\iint_{E\times E}
\mathbf{1}_{\{(\frac{\rho^{D\gamma}}{V^\gamma})^{\frac{p(1-s)}{1+D\gamma}}>\lambda
\}}V^{p\gamma-1}\Big]^{\frac{1-s}{1+D\gamma}}\,d\lambda\\
&\quad\le \int_{0}^{M^{\frac{p(1-s)}{1+D\gamma}}}
\Big[\iint_{E\times
E}V(x,y)^{p\gamma-1}\,d\mu(y)\,d\mu(x)\Big]^{\frac{1-s}{1+D\gamma}}\,d\lambda,
\end{split}
\end{equation}
where the last step used \eqref{eq:pvm}.
Note that, for any $x,y\in E$, $V(x,y)\le V(x_0,3R)$,
which, combined with \eqref{eq:ele} with $\alpha:=p\gamma$ and the
doubling condition of $\mu$, further implies that
\begin{align*}
\iint_{E\times E}V^{p\gamma-1}\le
\int_E \int_{\{y\in X:V(x,y)\le V(x_0,3R)\}}V(x,y)^{p\gamma-1}\,d\mu(y)\,d\mu(x)
\lesssim V(x_0,R)^{p\gamma+1}.
\end{align*}
This, together with \eqref{eq:lot1}, yields
\begin{align*}
\BNorm{\mathbf{1}_{E\times
E}\Big(\frac{\rho^{D\gamma}}{V^\gamma}\Big)^{\frac{p(1-s)}{1+D\gamma}}}{
L^{\frac{1+D\gamma}{1-s},1} (V^{p\gamma-1})}\le
M^{\frac{p(1-s)}{1+D\gamma}}\Big[cV(x_0,R)^{p\gamma+1}\Big]^{\frac{1-s}{1+D\gamma}},
\end{align*}
which clearly converges to $1$ as $s\nearrow 1$.
Substituting this estimate into \eqref{eq:BSVY-v2}, we obtain
\begin{equation*}
  \limsup_{s\nearrow 1}(1-s)\iint_{E\times E}
     \frac{|F|^p}{\rho^{ps}V}
   \leq(1+D\gamma)\BNorm{\frac{F}{\rho V^\gamma}}{
L^{p,\infty} (V^{p\gamma-1})}^p,
\end{equation*}
which completes the proof in this case.

\emph{Case 2: $\gamma\in(-\infty,0)$}.
Given $s\in(0,1)$, applying \eqref{eq:BSVY4-2} with
$q:=\frac{2}{1+s}\in(1,\infty)$,
we conclude that
\begin{equation}
\begin{split}
&\iint_{E\times E}\frac{|F|^p}{\rho^{ps}V}%\\
=\iint_{X\times X}\Big(\frac{|F|}{\rho V^\gamma}\Big)^{\frac{p(1+s)}{2}}
\mathbf{1}_{E\times E}\Big(\frac{|F| \rho}{V^\gamma}\Big)^{\frac{p(1-s)}{2}}
V^{p\gamma-1}\\
&\quad\le \frac{2}{1-s}
\BNorm{\Big(\frac{|F|}{\rho V^\gamma}\Big)^{\frac{p(1+s)}{2}}}{
L^{\frac{2}{1+s},\infty} (V^{p\gamma-1})}%\\
\BNorm{\mathbf{1}_{E\times E}\Big(\frac{|F| \rho}{V^\gamma}\Big)^{\frac{p(1-s)}{2}}}{
L^{\frac{2}{1-s},1} (V^{p\gamma-1})}\\
&\quad\le \frac{2}{1-s}
\BNorm{\frac{F}{\rho V^\gamma}}{
L^{p,\infty} (V^{p\gamma-1})}^{\frac{p(1+s)}{2}}\Norm{F}{L^\infty(\mu\times\mu)}
^{\frac{p(1-s)}{2}}%\\
\BNorm{\mathbf{1}_{E\times E}\Big(\frac{\rho}{V^\gamma}\Big)^{\frac{p(1-s)}{2}}}{
L^{\frac{2}{1-s},1} (V^{p\gamma-1})}.
\end{split}
\end{equation}
Since $\gamma\in(-\infty,0)$, for all $x,y\in E\subset B(x_0,R)$,
it follows from the doubling condition that
\begin{align*}
\frac{\rho(x,y)}{V(x,y)^\gamma}
\lesssim RV_0^{-\frac1D}V(x,y)^{\frac1D-\gamma},
\qquad
V_0:=V(x_0,R),
\end{align*}
where the implicit positive constants depend only on the doubling constant and $\gamma$.
Denoting
\begin{equation*}
  t:=\Big(\frac 1D-\gamma\Big)\frac{p(1-s)}{2}=\frac{p(1-\gamma D)(1-s)}{2D}>0,
\end{equation*}
it follows that
\begin{align*}
  &\BNorm{\mathbf{1}_{E\times E}\Big(\frac{\rho}{V^\gamma}\Big)^{\frac{p(1-s)}{2}}}{
  L^{\frac{2}{1-s},1} (V^{p\gamma-1})}
  \leq (cRV_0^{-\frac{1}{D}})^{\frac{p(1-s)}{2}}
  \BNorm{\mathbf{1}_{E\times E} V^t}{L^{\frac{2}{1-s},1}(V^{p\gamma-1})} \\
  &\quad=(cRV_0^{-\frac{1}{D}})^{\frac{p(1-s)}{2}}\int_{0}^{(cV_0)^t}
  \Big(\iint_{E\times E}\mathbf{1}_{\{V>\lambda^{\frac 1t}\}}
  V^{p\gamma-1}\Big)^{\frac{1-s}{2}}\,d\lambda.
\end{align*}
Recalling that $\gamma\in(-\infty,0)$ and using \eqref{eq:ele}, it follows that
\begin{equation*}
\begin{split}
  \int_0^{(cV_0)^t} &\Big(\iint_{E\times E} \mathbf{1}_{\{V>\lambda^{\frac 1t}\}}
  V^{p\gamma-1}\Big)^{\frac{1-s}{2}}\,d\lambda   \\
  &\leq \int_0^{(cV_0)^t} \Big[\int_{B(x_0,R)}\int_{\{y:V(x,y)>\lambda^{\frac 1t}\}}
  V(x,y)^{p\gamma-1}d\mu(y)\,d\mu(x)
  \Big]^{\frac{1-s}{2}}\,d\lambda   \\
  &\leq \int_0^{(cV_0)^t} \big( cV_0
  \lambda^{\frac{p\gamma}{t}}\big)^{\frac{1-s}{2}}d\lambda
  = (cV_0)^{\frac{1-s}{2}} \int_0^{(cV_0)^t} \lambda^{\frac{\gamma D}{1-\gamma D}}d\lambda
  \\
  &=(cV_0)^{\frac{1-s}{2}}(1-\gamma D)(cV_0)^{\frac{t}{1-\gamma D}}
 =(1-\gamma D)(cV_0)^{\frac{1-s}{2}(1+\frac{p}{D})}.
\end{split}
\end{equation*}
Substituting back, it follows that
\begin{align*}
\limsup_{s\nearrow 1}(1-s)\iint_{E\times E}
     \frac{|F|^p}{\rho^{ps}V}
   \leq2(1-D\gamma)\BNorm{\frac{F}{\rho V^\gamma}}{
L^{p,\infty} (V^{p\gamma-1})}^p,
\end{align*}
since all other factors above are some finite quantities raised to the power $1-s\to 0$.
This completes the proof of Lemma \ref{lem:BSVY105}.
\end{proof}

\end{document}
